\documentclass[11pt]{article}

\usepackage[margin=1in]{geometry}
\usepackage[T1]{fontenc}
\usepackage{lmodern}
\usepackage{microtype}
\usepackage{amsmath,amssymb,mathtools}
\usepackage{physics}
\usepackage{bm}
\usepackage{booktabs}
\usepackage{enumitem}
\usepackage{xcolor}
\usepackage[colorlinks=true,linkcolor=blue!50!black,urlcolor=blue!50!black]{hyperref}

\setlist{nosep,leftmargin=*}
\setlength{\parskip}{0.4em}
\setlength{\parindent}{0pt}

\newcommand{\gv}{\,|\,}
\newcommand{\Lop}{\mathcal{L}}
\newcommand{\cU}{\mathcal{U}}

\title{Answers to the questions in \texttt{FPE\_ch5.tex}}
\author{Presenter notes for the Chapter 5 talk}
\date{\today}

\begin{document}
\maketitle

\noindent
One section per \texttt{\% Question} comment in the slides, in slide order, plus the questions inside some
\texttt{\% TODO} comments and a section on symmetries. Notation is the deck's: $\partial_tp=-\partial_x(Ap)+\frac12
\partial_x^2(Bp)$, $B=D$ for constant noise (so $\langle\Delta x^2\rangle=Dt$), $\cU$ the physical potential and
$U=\cU/\gamma$. Equation numbers marked G2 are Gardiner's 2nd edition.

\tableofcontents

% =============================================================================
\section{It\^o, Stratonovich, and manifolds}
\emph{Slide: The Corresponding SDE.}

\paragraph{Where the two come from.} An SDE is shorthand for an integral, and $\int b(x)\,\dd W$ needs a rule for
where in each step $b$ is evaluated. It\^o evaluates it at the start of the step. Then $b$ is independent of the
increment, so $\langle b\,\dd W\rangle=0$. Stratonovich evaluates it at the midpoint. Expanding
$b(x_{\rm mid})\approx b(x)+\frac12b'(x)\,\Delta x$ and using $\Delta x\approx b\,\Delta W$ gives
\[
  b(x_{\rm mid})\,\Delta W \approx b\,\Delta W + \tfrac12\,b\,b'\,(\Delta W)^2 \to b\,\dd W + \tfrac12\,b\,b'\,\dd t .
\]
So the Stratonovich SDE $\dd x=a\,\dd t+b\circ\dd W$ is the It\^o SDE with drift $a+\frac12bb'$. With $b=\sqrt B$,
$\frac12bb'=\frac14B'$, which is the slide's $A=a+\frac14\partial_xB$. The FPE is always written with the It\^o drift.
In Stratonovich variables it has a neat ``sum of squares'' form:
\[
  \partial_tp = -\partial_x(a\,p) + \tfrac12\,\partial_x\bigl[b\,\partial_x(b\,p)\bigr].
\]
(Expand the last term to check that it equals $-\partial_x(\frac12bb'p)+\frac12\partial_x^2(b^2p)$.)

\paragraph{Which one is physical.} It depends on where the noise comes from.
\begin{itemize}
  \item Noise that is really smooth, with a short but nonzero correlation time, gives Stratonovich in the white-noise
        limit (the Wong--Zakai theorem). Ordinary calculus holds for smooth noise, and Stratonovich is the
        calculus that keeps the ordinary chain rule.
  \item Noise that acts in discrete kicks, with the state fixed before each kick, gives It\^o. Population models
        and finance are like this.
  \item A thermal particle whose friction depends on position, $\gamma(x)$, is neither. The overdamped limit of the
        Langevin equation must reproduce the Boltzmann density, and that requires the It\^o drift
        $A=-\cU'/\gamma+\frac12B'$, with $B=2k_BT/\gamma(x)$. This is sometimes called the isothermal or
        H\"anggi--Klimontovich convention. The interface section of the notes compares all three ($\chi=0$, $\frac12$, $1$).
\end{itemize}
When $B$ is constant all three agree, which is every example in the talk.

\paragraph{Manifolds: yes, this is the connection.} Under a change of variable $y=\varphi(x)$, It\^o's lemma gives
\[
  A_y=\varphi'A+\tfrac12\varphi''B,\qquad B_y=\varphi'^2B .
\]
$B$ transforms as a tensor (two upper indices), but the It\^o drift picks up a term with a second derivative of the
coordinate change, so $A^i$ is \emph{not} a vector field. The Stratonovich chain rule has no such term: $\dd
y=\varphi'(x)\circ\dd x$, so $a$ and $b$ both transform like components of vector fields. That is why SDEs on a
manifold are written in Stratonovich form,
\[
  \dd X=V_0(X)\,\dd t+\sum_\alpha V_\alpha(X)\circ\dd W^\alpha,
\]
with $V_0,V_1,\dots$ vector fields. The backward operator (the generator) is then coordinate-free,
$\Lop^\dagger=V_0+\frac12\sum_\alpha V_\alpha V_\alpha$, a sum of squares of first-order operators (H\"ormander's form).
Its adjoint with respect to a volume form $\mu$ is the forward operator,
$\Lop\rho=-\operatorname{div}_\mu(V_0\rho)+\frac12\sum_\alpha\operatorname{div}_\mu\bigl(V_\alpha\operatorname{div}_\mu(V_\alpha\rho)\bigr)$.

Example: Brownian motion on a Riemannian manifold has generator $\frac12\Delta_g$ (Laplace--Beltrami). In coordinates,
$\frac12\Delta_gf=\frac12g^{ij}\partial_i\partial_jf-\frac12g^{jk}\Gamma^i_{jk}\partial_if$, so its It\^o drift is
$-\frac12g^{jk}\Gamma^i_{jk}$: a drift with no force behind it, purely from the coordinates. The radial part of
free motion in $d$ dimensions, with drift $(d-1)D/2r$, is the simplest case.

\paragraph{What this means for the index notation on the slides.} Writing $A^i$ and $B^{ij}$ with upper indices is fine
in Cartesian coordinates, which is all Gardiner uses. Under nonlinear coordinate changes $B^{ij}$ is a tensor and $A^i$
is not. But the current $J^i=A^ip-\frac12\partial_j(B^{ij}p)$ \emph{is} a vector density (with $p$ a density), as it
must be if its flux through a surface is physical (slide ``Flux across any surface''). In one dimension, with
$p_y=p/\varphi'$, a short calculation gives $J_y=J_x$: the extra drift term cancels the extra derivative term.

% =============================================================================
\section{Why the overdamped limit is \texorpdfstring{$m\,\dd v\approx0$}{m dv = 0}}
\emph{Slide: The overdamped limit.}

Write the velocity equation as relaxation toward the local terminal velocity $v_*(x)=-\cU'(x)/\gamma$:
\[
  \dd v = -\frac{1}{\tau_v}\bigl(v-v_*(x)\bigr)\,\dd t + \frac{\sqrt{2\gamma k_BT}}{m}\,\dd W,\qquad \tau_v=\frac m\gamma .
\]
Integrate the momentum equation over a step $\Delta t$ with $\tau_v\ll\Delta t\ll\tau_x$:
\[
  m\,\Delta v = -\gamma\,\Delta x - \cU'\,\Delta t + \sqrt{2\gamma k_BT}\,\Delta W .
\]
\begin{itemize}
  \item The left side is bounded: $v$ is a thermal velocity, so $m\,\Delta v\sim\sqrt{mk_BT}$ whatever $\Delta t$ is.
  \item The noise impulse grows like $\sqrt{2\gamma k_BT\,\Delta t}$, and the force and friction impulses grow like
        $\Delta t$.
  \item Their ratio is $m\,\Delta v/\text{noise}\sim\sqrt{\tau_v/\Delta t}\to0$.
\end{itemize}
So on any time step much longer than $\tau_v$, the inertial term is negligible next to the others, and the
remaining terms must balance: $\gamma\,\dd x=-\cU'\,\dd t+\sqrt{2\gamma k_BT}\,\dd W$. ``$m\,\dd v\approx0$'' does
\emph{not} mean that $v$ is constant or small. $v$ fluctuates wildly, and in the limit it is white noise; what is small
is the momentum change compared with the impulses over a coarse step.

The upper limit $\Delta t\ll\tau_x$ is needed so that $\cU'$ is roughly constant over the step; the window
$\tau_v\ll\Delta t\ll\tau_x$ exists only if $\varepsilon=\tau_v/\tau_x=m|\cU''|/\gamma^2\ll1$.

\paragraph{A check you can quote.} For a free particle ($\cU=0$) the Langevin equation is solvable (Ornstein, 1917):
\[
  \langle\Delta x^2\rangle = \frac{2k_BT}{\gamma}\Bigl[t-\tau_v\bigl(1-e^{-t/\tau_v}\bigr)\Bigr]
  \;\to\;
  \begin{cases} (k_BT/m)\,t^2 & t\ll\tau_v \ \text{(ballistic)}\\ (2k_BT/\gamma)\,t=Dt & t\gg\tau_v\ \text{(diffusive)} \end{cases}
\]
The FPE describes only the second regime; that is the ``valid for $t\gg\tau_v$'' on the slide.

\paragraph{A caution.} With constant $\gamma$ the naive step is safe. If $\gamma$ depends on $x$, simply dropping $m\,\dd v$ and
reading the result as It\^o gives the wrong stationary state; the correct limit has the extra drift $\frac12B'$
(section 1).

% =============================================================================
\section{Smoluchowski's equation in Gardiner}
\emph{Slide: Smoluchowski's equation.}

Yes, in two places (2nd edition; the 4th renumbers, so use the index under ``Smoluchowski equation'' or
``adiabatic elimination''):
\begin{itemize}
  \item \S5.1 says some authors call the FPE the Smoluchowski equation, after Smoluchowski's early use of it.
  \item \S6.4, ``Adiabatic Elimination of Fast Variables'', starts from Kramers' equation (G2 \S5.3.6a) and derives the
        Smoluchowski equation as the leading order in $1/\gamma$ (G2 (6.4.18)). It then computes the next order, the
        corrected Smoluchowski equation (G2 (6.4.108)), which multiplies the diffusion by $1+m\cU''/\gamma^2$. That is
        the $1\pm\varepsilon$ on the slide.
\end{itemize}
So your ``supposedly'' is right: what the slide does by hand is the leading order of Gardiner's systematic
expansion. The slide now says so without the question mark.

% =============================================================================
\section{The Einstein relation}
\emph{Slide: Smoluchowski's equation.}

In physics units, with $D_{\rm phys}$ defined by $\partial_tp=D_{\rm phys}\,\partial_x^2p$ (so $\langle\Delta
x^2\rangle=2D_{\rm phys}t$),
\[
  D_{\rm phys}=\frac{k_BT}{\gamma}=\mu\,k_BT,\qquad \mu=\frac1\gamma\ \text{(mobility: velocity per unit force)}.
\]
It ties the size of the fluctuations (diffusion) to the dissipation (friction): the simplest fluctuation--dissipation
relation. One-line derivation: in a potential, zero current
$J=-(\cU'/\gamma)p-D_{\rm phys}\,p'=0$ gives $p\propto e^{-\cU/\gamma D_{\rm phys}}$, and this is the Boltzmann factor
$e^{-\cU/k_BT}$ only if $\gamma D_{\rm phys}=k_BT$.

History: Einstein (1905), Sutherland (1905) and Smoluchowski (1906). With Stokes friction $\gamma=6\pi\eta a$ it
becomes the Stokes--Einstein relation $D_{\rm phys}=k_BT/6\pi\eta a$, which Perrin used to measure Avogadro's number. In
the deck's convention $B=D=2D_{\rm phys}$, hence ``$D/2=k_BT/\gamma$''.

% =============================================================================
\section{Why the continuity equation fixes \texorpdfstring{$\bm J$}{J} only up to a divergence-free field}
\emph{Slide: Flux across any surface.}

The continuity equation $\partial_tp+\partial_iJ^i=0$ involves $\bm J$ only through its divergence. If $\bm J$
satisfies it, so does $\bm J+\bm K$ for any $\bm K$ with $\partial_iK^i=0$; in three dimensions, any $\bm K=\nabla\times\bm
M$. The flux of $\bm K$ through a \emph{closed} surface is zero (divergence theorem), so the total probability in every
region comes out the same. But the flux of $\bm K$ through an \emph{open} surface, such as $S_{12}$, need not be
zero. So the continuity equation alone cannot say how much probability crosses $S_{12}$. Gardiner's (5.2.12) computes
that crossing directly from two-time probabilities and finds that the particular $\bm J$ of (5.2.8) gives it. That is the
sense in which it is ``the physical'' current.

The one-dimensional version is the clearest. Divergence-free means constant, so the continuity equation fixes $J$ only
up to an added constant. On an interval the walls fix the constant. On a circle they cannot: the free circle and the
circle with constant drift $v$ both have the uniform stationary density $p_s=1/L$, so $p_s$ alone cannot tell you
whether probability circulates. The formula $J=vp-\frac12Dp'$ can: $J_s=v/L$.

% =============================================================================
\section{The walls figure: axis labels and \texorpdfstring{$L=2x_0$}{L = 2x0}}
\emph{Slide: Reflecting, absorbing, and periodic. (A \texttt{TODO} question.)}

The labels are consistent, but one fact was missing from the figure itself. Lengths are in units of the start $x_0$,
and time in units of $x_0^2/D$. All three panels use the \emph{same} Brownian path.
\begin{itemize}
  \item (a) Hard wall: the plotted path is $|x(t)|$. Reflected Brownian motion has the same law as $|W|$, which is
        the image solution $g(x-x_0)+g(x+x_0)$. The shaded band $x<0$ is the forbidden side.
  \item (b) Cliff: the path up to its first arrival at $0$ (the cross); dotted is what it would have done next.
  \item (c) Ring: yes, the circumference is $L=2x_0$, i.e.\ $L=2$ in these units. The ring coordinate is
        $x\bmod 2x_0$, which runs over $[0,2)$; the two horizontal lines at $0$ and $2$ are the same point of the ring
        (the cut). When the blue path reaches $2$ it reappears at $0$, and vice versa. Dotted grey is the lift, the
        unwrapped path, which is the original path again.
\end{itemize}
The slides now use a copy of the figure (made by \texttt{fpe-slides/scripts/slide\_figures.py} from the notes' own
plotting function) whose panel (c) title says ``ring of circumference $L=2x_0$''. The notes are unchanged.

% =============================================================================
\section{Winding on the circle: a quick overview}
\emph{Slide: Periodic boundaries, or: the circle.}

\paragraph{What it is.} A point on the circle is $x\bmod L$. Follow a path continuously and never jump back at the cut:
that gives the \emph{lift} $\tilde x(t)\in\mathbb R$, with $\tilde x(t)=x(t)+L\,N(t)$, where $N(t)\in\mathbb Z$ is the
net number of times around (anticlockwise minus clockwise). The winding is $\Delta\tilde x(t)=\tilde x(t)-\tilde x(0)$,
the net distance travelled around the ring. Panel (c) of the walls figure shows a path and its lift.

\paragraph{Why it is new.} The position on the circle tells you where the particle is; the winding tells you how it got
there. It exists only because the circle has a loop that cannot be shrunk to a point, so no interval has it. It is the
observable that measures \emph{transport} around the ring.

\paragraph{How it behaves.} At long times it drifts and spreads,
\[
  \langle\Delta\tilde x\rangle\simeq v_{\rm w}\,t,\qquad \operatorname{Var}\Delta\tilde x\simeq D_{\rm eff}\,t .
\]
\begin{itemize}
  \item The winding velocity is the stationary current times the circumference, $v_{\rm w}=LJ_s$. So $v_{\rm w}\neq0$
        exactly when $\Phi=\oint2A/B\neq0$, the non-equilibrium case.
  \item The effective diffusion constant is nonzero even in equilibrium. For $B=D$ and a periodic potential $U$, it is the
        Lifson--Jackson formula
        $D_{\rm eff}=D\big/\bigl(\langle e^{2U/D}\rangle\,\langle e^{-2U/D}\rangle\bigr)$, with $\langle\cdot\rangle$ the
        average over one period. A barrier makes it exponentially small: hopping between wells.
\end{itemize}
Examples: the free circle has $v_{\rm w}=0$ and $D_{\rm eff}=D$ (the ring changes where the particle is, not how far it
goes); constant drift has $v_{\rm w}=v$; the cosine well has $D_{\rm eff}=D/I_0(\kappa)^2$.

\paragraph{How it is computed.} Two equivalent ways.
\begin{itemize}
  \item First passage on the cover: the time $T_{\rm w}$ to first get a full circumference away, in either direction, is the
        exit time from $(x_0-L,x_0+L)$ on the line, an absorbing-interval problem. Without drift, windings are a random
        walk of steps $\pm L$, and $D_{\rm eff}=L^2/T_{\rm w}$.
  \item One eigenvalue: $\langle e^{s\Delta\tilde x}\rangle\sim e^{\Lambda(s)t}$, where $\Lambda(s)$ is the top
        eigenvalue of the forward operator with $\partial_x\to\partial_x-s$, acting on periodic functions. Then
        $v_{\rm w}=\Lambda'(0)$ and $D_{\rm eff}=\Lambda''(0)$. For constant drift, $\Lambda(s)=vs+\frac12Ds^2$.
\end{itemize}
The second way is Bloch's theorem in disguise: $\Lambda(s)$ is the lowest band of the periodic problem on the line,
continued to imaginary quasi-momentum $q=-is$. $D_{\rm eff}$ is the curvature of the band at $q=0$, the analogue of an
inverse effective mass.

\paragraph{Where it matters physically.} The dc voltage of a Josephson junction (the phase winds), a colloid driven
around a ring trap, rotary molecular motors, phase slips of an oscillator. For an interdisciplinary audience, the
one-liner is: \emph{a ring has a hole, and probability can circulate around it; the winding counts the circulation.}

% =============================================================================
\section{The squared-radius figure, panel by panel}
\emph{Slide: Feller's Classification, example.}

\paragraph{The process.} Take an overdamped particle in an isotropic harmonic trap in $d$ dimensions: each component obeys
$\dd x_i=-kx_i\,\dd t+\sqrt D\,\dd W_i$. Its squared distance from the centre, $y=\sum_ix_i^2$, obeys (It\^o's lemma,
$\dd x_i^2=D\,\dd t$)
\[
  \dd y=\bigl(dD-2ky\bigr)\,\dd t+2\sqrt D\sum_ix_i\,\dd W_i=\bigl(dD-2ky\bigr)\,\dd t+\sqrt{4Dy}\;\dd W .
\]
The last equation makes sense for any real $d$, which is how it also describes a population with demographic noise
($d$ is an immigration rate; $d=0$ is Feller's branching diffusion) and the CIR interest-rate model. Units:
$\ell^2=D/k$ for $y$, and $1/k$ for time. The noise $B=4Dy$ vanishes at $y=0$, so Feller's test decides what happens
there: exit for $d\le0$, regular for $0<d<2$, entrance for $d\ge2$.

\paragraph{(a) Sample paths, all from $y_0=\ell^2$.}
\begin{itemize}
  \item Blue, $d=3$: $|\bm x|^2$ of a 3D trapped particle. It never reaches $y=0$; a Brownian particle in two or more
        dimensions does not hit a given point. That is an entrance: unreachable from inside, but a particle started
        exactly at $0$ would leave at once.
  \item Green, $d=1$: $x^2$ of a 1D trapped particle. The black cross marks the first time it reaches $y=0$ (the moment $x$
        changes sign). It then leaves again: a regular boundary, reached and left. Because $y=x^2$, this path is the
        \emph{reflected} choice; the absorbed choice would stop at the cross.
  \item Orange, $d=0$: Feller's branching diffusion. At $y=0$ both drift and noise vanish, so once it arrives it stays
        (the thick orange line along the axis from $kt\approx1.5$): an exit, here extinction of the population.
\end{itemize}

\paragraph{(b) The spectrum as a function of $d$.} Two families of eigenvalues, in units of $k$:
\begin{itemize}
  \item blue, horizontal: reflected at $0$, $\lambda_n=2nk$, including the stationary $\lambda_0=0$. They exist for
        $d>0$, where $p_s\propto y^{d/2-1}e^{-y/\ell^2}$ is integrable at $0$;
  \item orange, falling: absorbed at $0$, $\lambda_n=(2n+2-d)k$. They exist for $d<2$, where the origin can be
        reached. The lowest, $(2-d)k$, is the decay rate of the survival probability.
\end{itemize}
The shading is Feller's classification read off the spectrum. In the exit window only the absorbed family exists:
nothing to choose. In the entrance window only the reflected family: the origin is never reached, so there is nothing to
absorb, and the absorption rate $(2-d)k$ has fallen to $0$ at $d=2$. In the regular window both exist, and choosing the
boundary condition chooses the spectrum. A consistency check: at $d=1$ the reflected rates $0,2k,4k,\dots$ are the
\emph{even} OU levels and the absorbed rates $k,3k,5k,\dots$ the \emph{odd} ones. That is the reflection symmetry of
section~15.

\paragraph{(c) The density near the boundary.} $\ell^2p(y,t\gv y_0)$ at $kt=0.25$, on log--log axes, to show the power
laws as $y\to0$:
\begin{itemize}
  \item reflected: $p\propto y^{d/2-1}$. For $d=3$ it vanishes like $y^{1/2}$ (blue); for $d=1$ it diverges like
        $y^{-1/2}$ (solid green), which is just the Jacobian of $y=x^2$: $p_y=p_x/2\sqrt y$ with $p_x$ finite at $x=0$;
  \item absorbed: $p$ tends to a nonzero constant (dashed green, $d=1$; dashed orange, $d=-1$).
\end{itemize}
The lesson: where $B$ vanishes, ``absorbing'' does not mean $p=0$, contrary to Gardiner's (5.2.14), which assumes $B>0$.
What vanishes is $p/p_s$, the backward function. In the $x$ variable the absorbed density is $2\sqrt y\,p_y\propto|x|$,
which does vanish at the origin, as an absorbing wall with $B>0$ should.

% =============================================================================
\section{Adjoints: do they need an inner product, and when are \texorpdfstring{$\Lop$ and $\Lop^\dagger$}{L and L-dagger} adjoint?}
\emph{Slides: It's Linear Algebra Time. (A \texttt{TODO} question and a \texttt{Question}.)}

\paragraph{No inner product is needed.} The general notion is the Banach-space adjoint (or transpose). Given a space $X$ and
a bilinear pairing $\langle f,p\rangle$ between $X$ and a second space $Y$ (typically its dual), the adjoint of
$T:X\to X$ is the operator $T^\dagger$ on $Y$ with $\langle f,Tp\rangle=\langle T^\dagger f,p\rangle$ for all $p$ and $f$.
The Hilbert-space adjoint is the special case $Y=X$, with the inner product as the pairing.

Here the natural pairing is the expectation value,
\[
  \langle f,p\rangle=\int f(x)\,p(x)\,\dd x=\mathbb E_p[f],
\]
between densities, $p\in L^1$ (absolutely integrable), and bounded observables, $f\in L^\infty$. $L^\infty$ is the dual of
$L^1$, so this is exactly the setting of a space and its dual. Probabilists say the same thing: $\Lop^\dagger$ is the
\emph{generator}, acting on observables, and $\Lop$ is its dual, acting on measures. The dagger suggests a Hilbert
space, but none is involved until we want orthogonality (section 10).

\paragraph{Yes: adjoint exactly when the boundary term vanishes.} The integration by parts (Lagrange's identity)
\[
  \langle f,\Lop p\rangle-\langle\Lop^\dagger f,p\rangle=\Bigl[-fJ-\tfrac12Bp\,\partial_xf\Bigr]_a^b
\]
holds for any smooth $p$ and $f$, whatever their boundary values. The differential expressions are therefore always
\emph{formal} adjoints. But an operator is a differential expression \emph{together with its domain}: the functions it is
allowed to act on, which is where the boundary conditions live. $\Lop$ (with its domain) and $\Lop^\dagger$ (with
its domain) are adjoint exactly when the boundary term vanishes for every $p$ in the first domain and every $f$ in the
second. Given the forward boundary conditions, this \emph{determines} the backward ones; this is how Gardiner derives
them in \S5.2.4 (G2 (5.2.65)--(5.2.69)):
\begin{itemize}
  \item absorbing, $p=0$: the term is $-fJ$, and $J$ at the wall is not zero, so we need $f=0$;
  \item reflecting, $J=0$: the term is $-\frac12Bp\,\partial_xf$, and $p$ at the wall is not zero, so we need
        $\partial_xf=0$ (for $B\neq0$);
  \item circle: the end points are the same point and everything is periodic, so the terms at $a$ and $b$ cancel;
  \item infinity: $p$ and $J$ decay, and the terms vanish for bounded $f$ and $f'$;
  \item singular boundaries where $B\to0$ (Feller): the term can vanish by itself, and that is what ``no condition
        needed'' means for entrance and natural boundaries;
  \item a partially absorbing wall, $J=\kappa p$, pairs with a Robin condition on $f$, $\frac12B\,\partial_xf=-\kappa f$.
\end{itemize}
Mismatched domains (say, forward reflecting with backward absorbing) are not adjoint, and the biorthogonality
$\int P_\lambda Q_{\lambda'}=\delta_{\lambda\lambda'}$ would fail.

% =============================================================================
\section{What is \texorpdfstring{$L^2$}{L2}? The three spaces}
\emph{Slide: Biorthogonality.}

$L^2(a,b)$ is the space of square-integrable functions, $\int f^2\,\dd x<\infty$, with the inner product $\langle
f,g\rangle=\int fg\,\dd x$. It is a Hilbert space. ``$\Lop$ is symmetric in $L^2$'' means $\langle
g,\Lop f\rangle=\langle\Lop g,f\rangle$ for all $f$ and $g$ in its domain. A weighted $L^2(w)$ has $\langle
f,g\rangle_w=\int fg\,w\,\dd x$. Three spaces matter:

\begin{center}
\renewcommand{\arraystretch}{1.25}
\begin{tabular}{@{}llll@{}}
\toprule
Space & Contains & Operator & Used for \\
\midrule
$L^1$ & densities & $\Lop$ & probability: stationary states, QSDs \\
$L^\infty$ & bounded observables & $\Lop^\dagger$ & survival, exit times, splitting \\
$L^2(1/p_s)$ & $\int p^2/p_s<\infty$ & $\Lop$ (symmetric if $J_s=0$) & orthogonality, gaps \\
\bottomrule
\end{tabular}
\end{center}

\begin{itemize}
  \item In plain $L^2(\dd x)$, $\Lop$ is not symmetric unless $p_s$ is uniform. Its adjoint there is $\Lop^\dagger$, the same
        expression as before, since both pairings are $\int fg\,\dd x$. So $Q_\lambda$, the eigenfunctions of
        $\Lop^\dagger$, are the left eigenvectors of $\Lop$, and $P_\lambda$ are its right eigenvectors.
  \item In $L^2(1/p_s)$, with zero stationary current, $\Lop$ is symmetric:
        $\int g\,\Lop p/p_s\,\dd x=\int p\,\Lop g/p_s\,\dd x$. That gives real eigenvalues, orthogonal eigenfunctions
        ($\int P_\lambda P_{\lambda'}/p_s=\delta_{\lambda\lambda'}$, the third form on the slide), and completeness:
        the spectral theorem.
  \item The spaces are linked by two maps. $p\mapsto p/p_s$ takes $L^2(1/p_s)$ to $L^2(p_s)$ and turns $\Lop$ into
        $\Lop^\dagger$ (Gardiner's $p=p_sq$, G2 (5.2.71)--(5.2.72)). $p\mapsto p/\sqrt{p_s}$ takes it to plain $L^2$, where
        $\Lop$ becomes minus a Schr\"odinger operator, $H=-\frac12D\partial_x^2+V$ with $V=A^2/2D+A'/2$.
  \item The weighted norm has a meaning: $\int(p-p_s)^2/p_s\,\dd x$ is the $\chi^2$ distance from equilibrium, and it
        only decreases,
        $\dv{t}\int p^2/p_s\,\dd x=-\int Bp_s\bigl[\partial_x(p/p_s)\bigr]^2\dd x\le0$. That is why all $\lambda\ge0$.
\end{itemize}
On a bounded domain all three spaces give the same discrete spectrum. On unbounded domains they can differ; the
continua of QSDs on the ``caution'' slide are positive eigenfunctions in $L^1$ at $\lambda$ that are not in the $L^2$
spectrum.

% =============================================================================
\section{The propagator line of the quantum-mechanics table}
\emph{Slide: Analogy with Quantum Mechanics.}

In quantum mechanics the propagator is the amplitude to go from $x_0$ to $x$ in time $t$,
\[
  K(x,t;x_0)=\langle x|e^{-iHt/\hbar}|x_0\rangle=\sum_n\langle x|n\rangle\langle n|x_0\rangle\,e^{-iE_nt/\hbar}
  =\sum_n\phi_n(x)\,\phi_n^*(x_0)\,e^{-iE_nt/\hbar},
\]
obtained by inserting the resolution of the identity $\sum_n|n\rangle\langle n|=1$. Any wavefunction evolves by
$\psi(x,t)=\int K(x,t;x_0)\,\psi(x_0,0)\,\dd x_0$.

The FPE line is the same construction. $p(x,t\gv x_0,0)$ is the kernel of $e^{t\Lop}$, and the resolution of the identity
is the completeness relation $\sum_\lambda P_\lambda(x)\,Q_\lambda(x_0)=\delta(x-x_0)$ (the $t=0$ case). Any density evolves
by $p(x,t)=\int p(x,t\gv x_0,0)\,p(x_0,0)\,\dd x_0$.

The differences are instructive.
\begin{itemize}
  \item The ``bra'' is $\phi_n^*$ in quantum mechanics because $H$ is Hermitian: its left eigenvectors are the conjugates of
        its right ones. $\Lop$ is not symmetric in plain $L^2$, so its left eigenvectors are a different set, the
        $Q_\lambda$.
  \item $K$ is an amplitude; $p(x,t\gv x_0,0)$ is a probability density. It is positive, and $\int p\,\dd x=1$ because
        $Q_0=1$ and $\int P_\lambda\,\dd x=0$ for $\lambda\neq0$.
  \item Setting $t=-i\tau$ turns $K$ into the Euclidean (heat) kernel $\sum_n\phi_n(x)\phi_n^*(x_0)e^{-E_n\tau/\hbar}$.
        Through $p=\sqrt{p_s}\,\psi$ the FPE propagator is exactly that kernel for $H$, dressed with the ground state:
        $p(x,t\gv x_0,0)=\sqrt{p_s(x)/p_s(x_0)}\;\sum_\lambda\psi_\lambda(x)\psi_\lambda(x_0)e^{-\lambda t}$.
\end{itemize}

% =============================================================================
\section{The spectra figure, column by column}
\emph{Slide: Example spectra.}

Each column is one problem, with its eigenvalues in its own natural unit (written under it). Black ticks are discrete
eigenvalues; shaded bands are continua; the orange arrow is the gap that controls the long-time behaviour.
\begin{enumerate}
  \item Reflecting interval: $0,1,4,9,\dots$ in units of $\pi^2D/2L^2$. $\lambda=0$ is the stationary state; the arrow is
        the relaxation rate $\lambda_1$.
  \item Absorbing interval: $1,4,9,\dots$ No zero: nothing survives. The lowest level is the decay rate of $S$, and the
        arrow $\lambda_2-\lambda_1=3$ is the rate of approach to the QSD.
  \item Free half-line: a continuum down to $0$, no gap, so relaxation is algebraic ($t^{-1/2}$).
  \item Drift into a reflecting floor (sedimentation): the stationary $\lambda=0$, isolated below a continuum that
        starts at $v^2/2D=1/2\tau$. The gap is the edge of the continuum.
  \item OU on the line: $0,1,2,\dots$ in units of $k$. Gap $k$ on an infinite domain, because the trap confines.
  \item Harmonic well, reflecting at the bottom and absorbing at $L=2.5\ell$ (the ``escape from a trap'' example).
\end{enumerate}
\paragraph{The last column.} Right: the lowest eigenvalue is not zero. The wall at $2.5\ell$ is absorbing, so there is
no stationary state and no $\lambda=0$; every particle eventually escapes. But the lowest eigenvalue is tiny,
$\lambda_0\approx0.005\,k$, too small to see on the axis, so the grey arrow labels it. It is the Kramers escape rate:
$\frac{2k}{\sqrt\pi}\frac L\ell e^{-L^2/\ell^2}$ with $L/\ell=2.5$ gives about $0.005\,k$ after the first correction.
The next level is at about $2.2\,k$, so the orange arrow, $\lambda_1-\lambda_0$, is about $450$ times $\lambda_0$. That ratio is
the metastability: the particle equilibrates in the well (rate $\approx2.2k$) long before it escapes (rate
$0.005k$). Compare the reflecting half-line, $0,2k,4k,\dots$: a far-away cliff barely moves the upper levels but turns
the $0$ into an exponentially small escape rate. SDE Explorer, opened from the ``escape from a trap'' slide, prints
$\lambda_0=0.004954$ and $\lambda_1-\lambda_0=2.232$.

% =============================================================================
\section{Why \texorpdfstring{$\lambda_2-\lambda_1$}{lambda2 - lambda1} sets the approach to the QSD}
\emph{Slide: Quasi-Stationary Distributions, absorbing walls.}

Expand and divide:
\[
  \tilde p=\frac pS
  =\frac{c_1P_1e^{-\lambda_1t}+c_2P_2e^{-\lambda_2t}+\cdots}{c_1\bar P_1e^{-\lambda_1t}+c_2\bar P_2e^{-\lambda_2t}+\cdots}
  =\frac{P_1}{\bar P_1}\cdot\frac{1+\frac{c_2P_2}{c_1P_1}\,e^{-(\lambda_2-\lambda_1)t}+\cdots}{1+\frac{c_2\bar P_2}{c_1\bar P_1}\,e^{-(\lambda_2-\lambda_1)t}+\cdots},
\]
with $c_n=Q_n(x_0)$ and $\bar P_n=\int P_n\,\dd x$. Dividing by $S$ removes the common factor $e^{-\lambda_1t}$, the
overall loss of probability. What is left is how fast the second mode dies \emph{relative to the first}, and that is
$e^{-(\lambda_2-\lambda_1)t}$. So $\tilde p=\text{QSD}+O(e^{-(\lambda_2-\lambda_1)t})$.

It is the same reason the power method of linear algebra converges at the ratio of the top two eigenvalues: normalising at
each step discards the overall scale, and only the relative weights evolve. If the start is symmetric and $c_2=0$, the
rate is $\lambda_3-\lambda_1$ instead (section 15).

% =============================================================================
\section{Plot Lab and SDE Explorer in the talk}
\emph{Slide: QSD example. (A \texttt{TODO} question.)}

Yes, the QSD example is a very good spot. In SDE Explorer you watch sample paths get killed at the walls while the histogram
of survivors settles onto the QSD, and the exit rate $-\dot S/S$ climbs to $\lambda_1$ live. Plot Lab's ``free diffusion,
absorbing walls'' page shows the same exact curves with sliders for $x_0$, $L$ and $D$.

The slides now have small buttons at the top right of these frames:
\begin{itemize}
  \item Plot Lab only: the walls figure (\texttt{physical}), the squared-radius example (\texttt{sqrt}), the spectra
        (\texttt{spectra}), the reductions tree (\texttt{reductions}: each box of the map links to its page), walls by
        images (\texttt{free-summary});
  \item Plot Lab and SDE Explorer: the QSD example, the box, the ring, constant drift, OU, and escape from a trap.
\end{itemize}
Every Plot Lab page also has its own ``Open in SDE Explorer'' button. I checked two of the links (the absorbing interval in
Plot Lab, and the escape problem in SDE Explorer) with both dev servers running; they open the right process and walls.

For the talk:
\begin{itemize}
  \item Run \texttt{npm run dev} in \texttt{plot-lab/} and \texttt{sde-explorer/} beforehand (ports 5174 and 5173), or
        deploy the built apps side by side and change the two URLs in the macros \verb|\plotlab| and \verb|\sdeexplorer| at
        the top of the \texttt{.tex} file.
  \item Some PDF viewers ask before opening an external link; do a practice click in the viewer you will present with.
\end{itemize}

% =============================================================================
\section{Symmetries in the solutions}
\emph{New slides: Symmetries (three frames), and remarks in the examples.}

A symmetry of the problem (domain, walls, $A$ and $B$) is an operation that commutes with $\Lop$. Commuting operators can be
diagonalised together, so symmetries label the eigenfunctions, forbid some overlaps (selection rules), and let one solve a
big problem from a smaller one. Each one below shows up somewhere in the examples.

\subsection{Reflection}
\begin{itemize}
  \item \textbf{Parity.} If the domain is symmetric about $c$, $B$ is even and $A$ is odd, the reflection $R:x\mapsto2c-x$
        commutes with $\Lop$. Eigenfunctions are even or odd about $c$: $\cos$ and $\sin$ in the box, $H_n$ of definite parity
        for OU.
  \item \textbf{Selection rules.} A start at $c$, or any even initial density, has no overlap with odd modes; an even
        observable never sees them. The box started at its centre relaxes at $\lambda_2=4\lambda_1$; in OU, $\langle x\rangle$
        is a single exponential at $k$ (only $H_1$), while $\langle x^2\rangle$ relaxes at $2k$ (only $H_2$).
  \item \textbf{Sectors are walls.} An even density has $J=0$ at $c$; an odd one has $p=0$ there. So the even sector of a
        symmetric problem is the half problem with a reflecting wall at $c$, and the odd sector is the half problem with an
        absorbing wall. That is the method of images (a $+$ image or a $-$ image), and it is ``folding'' in the notes: the
        mixed interval unfolds into an absorbing interval of twice the length; the squared-radius spectrum at $d=1$ splits
        into even and odd OU levels.
  \item \textbf{The double well as a tunnelling doublet.} In the symmetrised (Schr\"odinger) form, the ground state
        $\sqrt{p_s}$ is even and the first excited state is odd, nearly $\pm\sqrt{p_s}$ in the two wells. In quantum
        mechanics this pair is the tunnelling doublet, split by $\Delta E\propto e^{-S/\hbar}$. Here the splitting is
        $\lambda_1=2r$, twice the Kramers hopping rate, and the odd backward eigenfunction $Q_1\approx2\pi_c-1$ is the
        committor. Kramers' formula is the FPE's version of a WKB tunnelling estimate.
\end{itemize}

\subsection{Translations}
\begin{itemize}
  \item \textbf{Fourier modes.} On the line with constant coefficients, translations commute with $\Lop$: eigenfunctions
        are $e^{ikx}$, and $p$ depends only on $x-x_0$.
  \item \textbf{The circle as a quotient.} The circle is the line modulo translations by $L$, so its propagator is the line's
        summed over all translates (the cover), every image with a $+$ sign. A mirror imposes a condition and can flip a
        sign; a translation imposes nothing, it is just another way of arriving at the same point.
  \item \textbf{Doublets.} With rotations and a reflection, each nonzero level of the free circle is doubly degenerate:
        $\cos q_mx$ and $\sin q_mx$ about any point, which are the reflecting and the absorbing intervals of length $L/2$.
        In the notes' cosine well the doublets survive (exactly) without full rotation symmetry; a $\cos2x$ term would split
        them.
  \item \textbf{Bloch.} With periodic coefficients on the line, translations by $L$ still commute with $\Lop$. Eigenfunctions
        are Bloch waves, and the winding statistics come from the lowest band near zero quasi-momentum (section 7):
        $D_{\rm eff}$ is its curvature.
\end{itemize}

\subsection{Boosts, gauges and scaling}
\begin{itemize}
  \item \textbf{Galilean boost.} Constant drift on the line is free motion in the frame $x-vt$. It fails with walls (they
        would move), and on the circle it is the only way to remove the drift, since the gauge factor is not periodic.
  \item \textbf{The gauge} $p=e^{\alpha(x-x_0)-t/2\tau}u$ is not a symmetry but a similarity transformation: it maps the
        drift problem onto the free one, shifting every rate by $1/2\tau$ and turning reflecting walls into Robin walls. It is
        the symmetrisation $p=\sqrt{p_s}\psi$ for constant $A$.
  \item \textbf{Diffusive scaling.} $x\to cx$, $t\to c^2t$ leaves the free equation unchanged. So the free propagator is
        $(Dt)^{-1/2}$ times a function of $(x-x_0)/\sqrt{Dt}$; the half-line survival $\operatorname{erf}(x_0/\sqrt{2Dt})$ depends
        only on $x_0^2/Dt$; any problem with one length $L$ has rates $\propto D/L^2$ and times $\propto L^2/D$. Drift and
        confinement break the scaling by adding their own scales, $D/|v|$ and $\sqrt{D/k}$.
  \item \textbf{OU's map} $y=xe^{kt}$, with the clock $\theta=(e^{2kt}-1)/2k$, is a time-dependent rescaling that turns OU
        into free diffusion. It keeps a wall at $x=0$ fixed (it is a fixed point of the rescaling) and moves any other wall,
        which is why the escape problem stops at Kummer functions.
\end{itemize}

\subsection{Time reversal and detailed balance}
\begin{itemize}
  \item A stationary process is \emph{reversible} if a film of it run backwards has the same statistics
        (detailed balance):
        \[ p(x,t\gv y,0)\,p_s(y)=p(y,t\gv x,0)\,p_s(x). \]
  \item In the expansion this holds exactly when $P_\lambda=p_sQ_\lambda$. Then
        \[ P_\lambda(x)\,Q_\lambda(y)\,p_s(y)=p_s(x)\,Q_\lambda(x)\,Q_\lambda(y)\,p_s(y), \]
        which is symmetric in $x\leftrightarrow y$. So Gardiner's
        relation $P=p_sQ$ \emph{is} time-reversal symmetry. Equivalently, $J_s=0$; equivalently, $\Lop$ is symmetric in
        $L^2(1/p_s)$, so the spectrum is real.
  \item In one dimension with a wall (or $J(\pm\infty)=0$), $J_s$ is a constant that vanishes at the wall, so detailed
        balance is automatic. Gardiner says it is ``almost trivial in one variable systems''.
  \item Breaking it needs a loop or more dimensions. On a circle the obstruction is $\Phi=\oint2A/B$; a driven ring has
        $J_s\neq0$, complex-conjugate pairs of eigenvalues ($\lambda_m=\frac12Dq_m^2+ivq_m$ for constant drift), and a density
        that circulates as it relaxes. In $n$ dimensions the obstruction is a curl: detailed balance needs
        $B^{-1}(2\bm A-\partial\cdot B)$ to be a gradient (Gardiner's potential conditions).
\end{itemize}

\subsection{Hidden symmetries}
\begin{itemize}
  \item \textbf{Forward--backward duality.} $\Lop$ and $\Lop^\dagger$ have the same spectrum, with $P_\lambda\leftrightarrow
        Q_\lambda$ and the wall conditions paired (section 9).
  \item \textbf{Partner potentials (supersymmetry).} With $B=D$, the symmetrised operator factorises,
        $H=\mathcal A^\dagger\mathcal A$ with $\mathcal A=\sqrt{D/2}\,(\partial_x-A/D)$, and $\mathcal A\sqrt{p_s}=0$. The
        partner $\mathcal A\mathcal A^\dagger$ is the same problem with $U\to-U$, and the two share every nonzero rate. For
        OU the partner of the trap is the inverted trap, whose rates are $k,2k,\dots$: OU's ladder with the $0$ removed. In
        the notes this explains the ``ladders by algebra'' box of the tree, and why $U\to-U$ swaps exit and entrance
        boundaries.
  \item \textbf{Coordinate invariance.} Feller's type of a boundary does not change under a change of variable
        ($y=r^2$, Lamperti, $\theta=2\arctan x$), because his integrals do not. A boundary can look regular in one chart and
        singular in another, but its type belongs to the process. Stratonovich form makes the same point for the SDE
        itself (section 1).
\end{itemize}

\end{document}
